\documentclass[a4paper,10pt]{article}

% include some standard latex packages
\usepackage{geometry}
\geometry{a4paper, top=25mm, left=25mm, right=25mm, bottom=30mm}

\usepackage{graphicx}
\usepackage{amssymb}
\usepackage{amsfonts,amsmath,bm}
\usepackage[mathscr]{eucal}
\usepackage[amssymb,thinqspace,thinspace]{SIunits}
\usepackage{hyphenat}
\usepackage{subfigure}
\usepackage{float}
\usepackage{helvet}
\renewcommand{\familydefault}{\sfdefault}

% include our own response package
\usepackage{layout_response_to_review}

%% Include your own defines located in 'defines.tex'
\include{defines}

%% Enable cross-references to main document of the paper
% Use either xr or xr-hyper and hyperref
%\usepackage{xr}
\usepackage{xr-hyper}
\usepackage{hyperref}
\externaldocument{main}

%% redefine thanking if necessary
%\renewcommand{\Thanks}{The authors would like to thank the reviewer ...}

%% Insert title and list of authors
\title{{\small Response to the review of the manuscript:}\\Trilinos: Enabling Scientific Computing across Diverse Hardware Architectures at Scale}
\author{Matthias Mayr \and Alexander Heinlein \and Christian A. Glusa \and Sivasankaran Rajamanickam \and Maarten Arnst \and Roscoe A. Bartlett \and Luc Berger-Vergiat \and Erik G. Boman \and Karen D. Devine \and Graham Harper \and Michael A. Heroux \and Mark Hoemmen \and Jonathan J. Hu \and Brian Kelley \and Drew P. Kouri \and Paul Kuberry \and Kyungjoo Kim \and Kim Liegois \and Curtis C. Ober \and Roger P. Pawlowski \and Carl Pearson \and Mauro Perego \and Eric T. Phipps \and Denis Ridzal \and Nathan V. Roberts \and Christopher M. Siefert \and Heidi K. Thornquist \and Romin Tomasetti \and Christian R. Trott \and Raymond S. Tuminaro \and James M. Willenbring \and Michael Wolf \and Ichitaro Yamazaki}
\date{\today}

\begin{document}

\maketitle

\ThanksToAll

%%%%%%%%%%%%%%%%%%%%%%%%%%%%%%%%%%%%%%%%%%%%%%%%%%%%%%%
%% Reviewer 1
%%%%%%%%%%%%%%%%%%%%%%%%%%%%%%%%%%%%%%%%%%%%%%%%%%%%%%%
\pagebreak
\Review{1}

%%%%%%%%
\Comment{Comments:
The paper presents high-level details of the Trilinos open-source library.  Each package in Trilinos is summarized.
While no results are shown, having a general description of the library is beneficial to the community.
I recommend the paper be published after a few minor changes.}

\Comment{Minor changes to paper:
\begin{itemize}
\item Change the section 2.4 heading for Tpetra to be "Distributed-Memory: Tpetra", and add several paragraphs at the start of the Tpetra library describing its general utility to all codes that seek MPI+Kokkos parallelism, adding references to examples that use Tpetra for HPC applications.
Reason:
The Tpetra library utility extends way beyond linear solvers.  The authors should highlight its full use, starting with its generality and then discuss specific applications of the library like using it with linear solvers.  The discussion of Tpetra library should start with how it supports distributed data types and parallel algorithms to leverage MPI+Kokkos, enabling users to implement multi-physics codes that exploit modern HPC machines.  One such example is the LANL open-source Fierro code (please cite it), which builds its MPI+Kokkos parallelism using Trilinos's Zoltan+Tpetra library (plus Kokkos kernels etc.).  In other words, code projects that may not need the Trilinos solvers (e.g., explicit hydrodynamic methods) will greatly benefit from using Zoltan+Tpetra (plus Kokkos etc).   Furthermore, by building on Tpetra, codes like Fierro can also seamlessly interface with the Trilinos solvers, allowing for example, novel topology optimization research leveraging MueLu and ROL, please see and cite:
\\\\
Diaz, A., Morgan, N. \& Bernardin, J. Parallel 3D topology optimization with multiple constraints and objectives. Optim Eng 25, 1531–1557 (2024). https://doi.org/10.1007/s11081-023-09852-6\\\\
Note, the above optimization work by Diaz et. al. only considered MPI parallelism; however, the authors have presented work at conferences on methods in Fierro that use MPI+Kokkos parallelism, leveraging Trilinos's Zoltan+Tpetra (etc) libraries.\\\\
Yes, Tpetra can be used with the linear solvers, but it's so much more than that.  It's a core library, not a linear algebra library.  As such, Tpetra is foremost a library that software teams should use when implementing MPI+Kokkos parallelism, impacting a huge community of researchers who may not need linear solvers.  Tpetra second application is for leveraging Trilinos solvers — linear and non-linear solvers and ROL for instance.
\end{itemize}~%
}

\todo{Matthias: The reviewer (probably Nathaniel Morgan) explicitly asks for citing a specific code and paper and provides an outline on how he wants the Tpetra section to be rewritten. I do not agree with his assessment, yet we need to be aware that this is his only request for change and we are about to deny it.}

\Response{
  Tpetra provides sparse distributed linear algebra.
  This does not mean that it can only be used in the context of the solution of linear systems.
  Trilinos is very well aware of Tpetra use cases besides solving linear systems:
  for example, finite element discretizations provided by Intrepid2 and Panzer rely on Tpetra data structures.
  Optimization packages such as ROL can leverage Tpetra objects.
  In our view, ``linear algebra'' does not equate to ``solution of linear systems of equations''.
}

\Comment{%
\begin{itemize}%
\item Additional comments on Trilinos, the authors do not need to alter the paper based on these comments.
\begin{itemize}
\item[A)] In addition to expanding the discussion in the paper on TPetra, the authors should update the web documentation for Trilinos calling attention to Tpetra's general utility for developers that are creating MPI+Kokkos codes.
\end{itemize}~%
\end{itemize}~%
}%

\Response{
  The Trilinos team is in the process of removing the old Epetra package stack.
  The presentation of Trilinos provided capabilities and their documentation should become a lot less confusing once everything is built around a single sparse linear algebra package.
}

\Comment{%
\begin{itemize}
\item[]
\begin{itemize}
\item[B)] Users of Trilinos would greatly benefit from faster compile times.  For example, when building with CUDA, it takes nearly a day to compile.
Even CPU builds of Trilinos take a long time.
The principal challenge with using Trilinos is long compile times.
There was no mention of compile times in the paper.
\end{itemize}~%
\end{itemize}~%
}

\Response{
  The compile times of Trilinos depend on the number of packages and backends enabled.
  Excessive compile times can occur, but are often caused by the underlying software stack.
  Trilinos tests every PR in its continuous integration, including CPU and CUDA builds.
  The turnaround for building and testing Trilinos is generally on the order of one hour, but admittedly relies on \texttt{ccache}.
  We will consider documenting best practices for achieving fast builds on the Trilinos web page or GitHub wiki.
}

\Comment{
\begin{itemize}
\item[]
\begin{itemize}
\item[C)]Please update the Trilinos anaconda distribution with each release.  The current anaconda version of Trilinos is very dated.  Also, it would be great to offer CUDA and HIP anaconda distributions of Trilinos.  Having anaconda distributions would help users adopt Trilinos, eliminating the need to compile it.  Also, the paper never mentions that Trilinos is available as an anaconda package, not sure it needs mentioning though.
\end{itemize}~%
\end{itemize}~%
}

\Response{
  The Trilinos package provided by Anaconda is based on PyTrilinos which uses the old Epetra stack.
  It will be archived in the fall and is therefore not mentioned in the paper.
  There are currently no pip or anaconda packages for its successor PyTrilinos2, but we will consider providing them in the future.
}

%\Response{Respond to the reviewer's comment.
%Explain why we did changes to the manuscript based on his/her comment.
%To quote a specific comment by the reviewer or a sentence from the manuscript,
%use the \texttt{\textbackslash Quote\{Let's quote the reviewer!\}} command
%yielding \QuoteReviewer{Let's quote the reviewer!}
%The special formatting of this quote will help to identify it as a quote.}
%
%\Changes{Quote the changes sentence/paragraph/figure here.}

%%%%%%%%%%%%%%%%%%%%%%%%%%%%%%%%%%%%%%%%%%%%%%%%%%%%%%%
%% Reviewer 2
%%%%%%%%%%%%%%%%%%%%%%%%%%%%%%%%%%%%%%%%%%%%%%%%%%%%%%%
\pagebreak
\Review{2}

%%%%%%%%
\Comment{The manuscript under review presents the quite well-known software package framework Trilinos.
The more than 30 pages document provide an update overview of the original presentation of Trilinos which dates back to 2005.
The paper delivers the update with a systematic overview of all the packages provided by Trilinos divided in a number of categories named ``product areas''
which are meant to give a logical division of the various packages included in Trilinos.\\\\
The manuscript provides a systematic description of the areas and the packages within. Each
package is described in its own subsection and connections or dependencies between packages within
the same areas are described in details. Each package is described using many abstract concepts
that are not always accessible by the average reader. For example, while a typical reader could be
an expert in sparse linear algebra constructs and concepts it may be completely foreign to concepts
in discretization or problem analysis.\\\\
The paper provides a lot of information and it definitely constitute a useful and overdue update
on the evolution of Trilinos in the last 20 years. Despite its many achievements, it feels sometimes
a bit too terse and fails to breach the communication barrier for non-expert in some of the specific
areas or packages described. This is an important shortcoming because it may fail to attract
additional users to Trilinos which, I believe, would be a desirable outcome of the manuscript.
This could be also the result of a manuscript which has probably seen many authors contributing.
Despite the fact that the language level is fairly uniform, there are subsections that are clearly
more effective in delivering their content than others (see detailed comments).}

\Response{We are aware that parts of the paper can leave the impression of being very terse and might require some background knowledge on the underlying mathematical methods.
Given the variety of algorithms included in Trilinos, it might be too big of a task to describe all algorithms in full detail in a single manuscript, especially given page limits of scientific journals.
As the reviewer observes, the manuscript already exceeds 30 pages.
The main goal of this manuscript is to provide an update on recent changes to the overall Trilinos project,
in particular its adoption of the Kokkos programming model to enable performance portability across hardware architectures.
To maintain a finite scope, the reader is referred to the packages' documentation or individual publications for details on concrete algorithms in specific software packages.}


\Comment{Overall, I believe the manuscript is definitely worth considering for publication as long as the
detailed comments listed below are taken into consideration.}

\noindent\textbf{Section 1}
\Comment{Page 3, subsec 1.1, paragraph 3 – ``... execution and memory spaces...'' Please provide an
explicit explanation of is intended with the spaces above.}

\Response{We have added short explanations for the purpose of execution and memory spaces.}

\Comment{Page 3, subsec 1.2, par. 2 – Please define explicitly the meaning of the term ``snapshotting''}

\Response{We have included the following explanation in line with the common use of this term in software development:}

\Changes{While most packages are developed natively in the Trilinos repository, some packages are developed externally and ``snapshotted'' into Trilinos,
i.e., incorporated as fixed copies of their source code at a specific point in time.}

\Comment{Page 5, par 1 – What is exactly intended with ``high-level'' algorithms? Is there something like
“low-level” algorithms?}

\Response{We used ``high-level'' to distinguish from lower level algorithms from sparse linear algebra and linear solvers.
We agree that this terminology lacks both rigor and generality.
Since the distinction is not needed in this context, we removed both expressions, ``high-level'' and ``low-level''.}

\Comment{Page 5, par 1 and 2 – automatic differentiation is mentioned in both the ``Non-linear Solvers
and Analysis Tools'' and ``Discretization Tools''. I would suggest that already at this stage an
explanation of the difference between these AD is provided}

\Response{AD plays an important role in many aspects of computational science.
Thus, it is used in different packages of Trilinos.
Section 4.3 introduces Sacado, Trilinos' AD framework, which is formally part of ``Nonlinear solvers and analysis tools''.
The discretization schemes provided by Trilinos use AD capabilities, which is why AD was mentioned in this paragraph as well.
We have rephrased the paragraph to make it clearer that these are not different types of AD,
but that AD used by discretization schemes refers back to the AD framework of ``Nonlinear solvers and analysis tools''.}

\Comment{Page 5, subsec 1.3 – I would strongly recommend to give in this section an overview of how
much has changed from the 2005 paper and current one. Implicitly, such a overview would give the
extent of how much the Trilinos project has evolved and the volume of the original contributions
included in this evolution.}

\Response{We have added the following paragraph for clarification:}

\Changes{In comparison to the first Trilinos publication from 2005~\cite{Heroux2005a}
and its updates from 2012~\cite{Heroux2012,HerouxIntro2012part1,HerouxIntro2012part2,pawlowski2012automating,Oldfield2012,bochev2012,Boman2012,Baker2012,Kokkos2012,pawlowski2012automatingpart2,Spotz2012,Long2012,Morris2012,Howle2012,Bavier2012a,Gaidamour2012},
the present manuscript not only provides updates on newly added capabilities,
but more importantly reflects the overhaul of Trilinos towards performance portability through the full integration of the Kokkos ecosystem~\cite{edwards2014,trott2021kokkos,trott2022kokkos}.
Many packages described in the original Trilinos publication~\cite{Heroux2005a} have been replaced by Kokkos-capable successor implementations,
among them (in alphabetical order)
Amesos2 replacing Amesos,
Belos replacing AztecOO,
Ifpack2 replacing Ifpack,
MueLu replacing ML,
and Tpetra replacing Epetra.
The present manuscript introduces these new packages and their capabilities and outlines their interplay with Kokkos.}

\noindent\textbf{Section 2}

\Comment{Page 6, subsec 2.2, par 4 -- ``batched dense and sparse linear algebra'' The reference here is on linear solvers only. What about BLAS? Also later the text says ``..on these general capabilities''; it is unclear what capabilities it refers to.}

\Response{\todo{respond}}

\Comment{Page 5-6 -- A general comment is in order here. It is not clear how much of the Kokkos development is carried out under the Trilinos framework and how much is external/independent. Could you please make this as much as possible explicit?}

\Response{Kokkos and its subprojects are nowadays developed independently and outside of Trilinos. For convenience, Kokkos releases are snapshotted into Trilinos. We have added a sentence to clarify this:}

\Changes{While the development of Kokkos Core has originated as a package of Trilinos, Kokkos has become an independent project in \todo{add year}
and is nowadays snapshotted into Trilinos for the pure convenience of dependency management.}

\Response{Furthermore:}

\Changes{Development of Kokkos Core itself is carried out in its own GitHub repository\footnote{https://github.com/kokkos/kokkos}.}

\Comment{Page 6, subsec 2.3 -- From the way Teuchos is described I would rather call it a ``programming interface''. The word ``tools'' convey the wrong impression, in my opinion.}

\Response{\todo{respond}}

\Comment{Page 6, subsec 2.4 -- ``dense vectors''. In section 2.2 it is claimed that Kokkos kernels can provide ``batched dense linear algebra'' but is this section Tpetra supports only dense vectors not dense matrices. I have hard time reconciling these the ability of doing dense linear algebra at the node level with the apparent inability of doing the same at the distributed level. A clarification would help here.}

\Response{
  A \texttt{MultiVector} is a collection of multiple vectors. It can therefore be used to represent a dense matrix.}

\Comment{Page 7, par 1 -- Could you give a simple example of ``Map data structure''.}

\Response{We forgo the inclusion of a detailed example as we believe that the included description matches the level of detail of the paper. Detailed examples of maps can be found in the Tpetra documentation.}

\Comment{Page 7, 1st bullet point -- Why associating BLAS-1 like kernels with TSQR?}

\Response{The bullet describes the \texttt{Vector} and \texttt{MultiVector} classes provided by Tpetra. \texttt{MultiVectors} are used in Trilinos to hold tall-and-skinny dense matrices. This explains why TSQR is listed under this bullet.}

\Comment{Page 7, 3rd bullet point -- ``Associated kernels...''. Can you be a bit more explicit in saying how kernels are associated with matrix storage format?}

\Response{The listed kernels need to be aware of the specific data storage format as they directly access data buffers. We have made this explicit in the text.}

\Comment{Page 7, subsec 2.5 -- ``.. the algorithms can be ... implementation such ..''. This whole sentence is unclear to me. When it talks about Krylov solver, is this a solver the Mat-Vec can be executed using kernels from another library? Is this the level of abstraction or something else?}

\Response{Yes, this provides an abstract layer.}

\Comment{Page 7, subsec 2.5, 1st bullet point -- ``... performance can be difficult.'' I find this paragraph and especially the sentence ending the words above confusing. What I understand is that Thyra and Tpetra have overlapping scopes. Tpetra targets distributed linear algebra but it looks it also provides methods like AXPY. What is then the distinction with Thyra?}

\Response{Thyra provides abstract interfaces. Tpetra is one of the possible concrete implementations of these interfaces.}

\Comment{Page 8, line 1 -- ``factories''. No idea what it is intended with this term.}

\Response{The following sentence describes what a factory is.}

\Comment{Page 8, par 1 -- ``... with the ANAs ... given in section 3.8'' Despite the text gives an example, I feel lost and find the example more confusing than illuminating.}

\Response{We have rephrased this sentence to clarify, that Stratimikos is a concrete implementation of this concept inside Trilinos:}

\Changes{[...] The solver interfaces allow one to wrap any concrete linear solver or preconditioner for use with the ANAs. As a concrete instantiation of this concept, Trilinos implements the Stratimikos package that wraps all of the Trilinos preconditioners and linear solvers in Thyra abstractions. [...]}

\Comment{Page 8, subsec 2.6, par 3 -- ``Zoltan2 ... (via Kokkos)''. Another source of confusion. I understood that what is described here is Tpetra’s goal}

\todo{Matthias: I don't understand the problem. The manuscript is very clear here. Should we just remove the line break between the second and third paragraph to reduce the room for misinterpretation by moving the statements about parallelism closer to the description of Zoltan2's capabilities?}

\Response{The first and second paragraph of subsection 2.6 clearly describe the purpose of Zoltan2: partitioning, load balancing, and other combinatorial algorithms.
While Tpetra offers distributed linear algebra, it does not compute the layout of the distribution itself.}

\noindent\textbf{Section 3}

\Comment{Page 9, line 2 -- I am not familiar with the ``traits mechanism'' concept. I would encourage to be less terse.}

\Response{We have slightly restructured the paragraph to better convey the fact that the following sentence gives examples of traits classes.}

\Comment{Page 9, par 1 -- ``... powerful solver managers..'' too abstract and not very operational description. I would suggest to clarify.}

\Response{\todo{respond}}

\Comment{Page 9, par 2 -- ``Belos .... solve the right-hand side'' An example on how the manager works across solvers would be very helpful here.}

\Response{\todo{respond}}

\Comment{Page 10, line 1 -- ``H(curl) and H(div) discretizations ...'' this sentence may result cryptic for the person unfamiliar with the notation.}

\Response{We have replaced H(curl) and H(div) discretizations with the types of physics that these discretizations are used for.}

\Comment{Page 10, subsec 3.3, par 1 -- ``wide range of challeging problems'' Like what? Perhaps the paragraph starting with ``FROSch has been applied...'' should be moved to an earlier position in the subsection.}

\Response{The sentence
\begin{quote}
  Besides \emph{parallel scalability}, FROSch emphasizes \emph{applicability and robustness across a wide range of challenging problems} while supporting \emph{algebraic construction} of the preconditioners, that is, solely from the fully assembled system matrix.
\end{quote}
highlights several important properties of FROSch. In the next paragraph, we first focus on the algorithmic aspects that enable parallel scalability and algebraicity. The sentence in question, which has also been slightly reformulated, then begins the following paragraph and addresses the wide range of applicability. We find it natural to separate these two aspects into distinct paragraphs and hope the referee agrees with this structure.}

\Comment{Page 10, subsec 3.3, last sentence -- ``..with performance gains... on GPU'' It would help to be more explicit here by at least mentioning the order of magnitude of the gains.}

\Response{We have added
\begin{quote}
  speedups of up to $4.8\times$ are observed when comparing ILU solve times on CPU and GPU
\end{quote}
to indicate best case performance gains from using GPUs.
}

\Comment{Page 11, par 1 -- ``MueLu is extensible and allow for the reseacrh and developemnt of new...'' The package allows for the research and development? Unclear what is meant by the author. Just two lines below ``Blue Gene/Q system'' This is a quite old system and the data referring to it is quite outdated. I suggest to cite something more recent.}

\Response{We have modified the statement in question to better explain how MueLu's flexibility can be leveraged for algorithmic development.}

\Comment{Page 12, subsec 3.5 -- The direct solvers described in this section seem to be application specific. Nothing wrong with it but it should be perhaps emphasized if they could be used with other problems.}

\Response{We have rearranged the content of this section, such that we first describe Amesos2 with its solvers with general applicability and then detail more application-targeted solvers in ShyLU.
We have added statements to make this distinction clear to the reader.}

\Comment{Page 12, subsec 3.5, par 2 -- I am not sure of the definition of BTF here. Is it really block triangular or should it be block tridiagonal? Please check. At the end of the paragraph the acronym KLU is undefined.}

\Response{\todo{respond}}

\Comment{Page 12, subsec 3.5, par 3 -- ``In addition... FROSch).'' So both Taco and Basker are node-level solvers only? This was not clear. Perhaps it could be made more explicit.}

\Response{\todo{respond}}

\Comment{Page 12, subsec 3.5, last par -- From this paragraph I evince that Amesos2 is only an interface not a solver package. Is this correct? Please clarify.}

\Response{Amesos2 implements one direct solver, KLU2, itself and provides interfaces to direct solvers in third-party libraries as well as in the Trilinos package ShyLU.
We have clarified that in the description of the Amesos2 package.}

\Comment{Page 12-13, subsec 3.6 -- The whole subsection is a bit unclear. I don’t have a clear idea of what Teko does. An example could perhaps clarify its scope.}

\Response{We have added a statement to further describe the purpose of the Teko package.}

\Comment{Page 13, subsec. 3.8 -- The description in this subsection sounds very nice indeed. I would very much like an example of the versatility of the interface. Also, why Anasazi is not included? Is the requirement of having a RHS such a constraint for the interface to exclude eigensolvers?}

\Response{The Stratimikos package was designed with the solution of linear systems \(Ax=b\) in mind, allowing the to access a wide variety of linear solvers and preconditioners in a uniform fashion. While it is true that it could be extended to allow for eigenvalue problems \(Ax=\lambda x\), Anasazi is the only eigensolver package in Trilinos. On the other hand, Anasazi itself implements Thyra and Tpetra interfaces.}

\noindent\textbf{Section 4}

\Comment{Page 14, subsec 4.1, par 1 -- ``custom implementations'' of what? Please explain.}

\Response{NOX algorithms are built on abstract base classes for solvers, directions, line searches and stopping criteria. Users can write derived clases for each of these base classes and register them with nox. This allows users to fully customize the algorithm. The statement has bee expanded on to clarify.}

\Comment{Page 14, subsec 4.2, par 1 -- ``any computational hardware'' A bit too eager of a statement. Does it include neuromorphic architectures of FPGAS or Quantum accelerators? I doubt it. Just tone it down. Also later ``many-core systems'' is a term that refers already to accelerators as opposed to multi-core systems. Please correct.}

\Response{\todo{respond}}

\Comment{Page 14, subsec 4.2, first bullet point -- ``ROL::Vector class...'' Can you explain why there is the need for another abstraction for vectors on top of the one already provided by other packages?}

\Response{\todo{respond}}

\Comment{Page 15, subsec 4.2, 2nd bullet point -- ``risck-averse problems'' I think there should be a definition of this concept for the non-expert reader.}

\Response{\todo{respond}}

\Comment{Page 15, subsec 4.2, 3rd bullet point -- ``data science and machine learning'' Perhaps a reference here? LASSO? Else?}

\Response{\todo{respond}}

\Comment{Page 15, subsec 4.2, last bullet point -- I find the description of this package based on the assumption that the only viable discretization is based on finite elements. What about finite differences? Perhaps this should be exemplified.}

\Response{\todo{respond}}

\Comment{Page 16, subsec 4.4, par 2 -- ``... to be outer-spatial and inner stochastic'' I have no idea what this sentence mean. This is a bit too cryptic.}

\Response{We revised this paragraph and included the following sentences to better clarify what is meant by outer-spatial and inner-stochastic:}

\Changes{In particular, stochastic Galerkin methods applied to PDEs employ a tensor product discretization between the spatial and stochastic discretizations that generates a two-level sparsity pattern in the resulting stochastic Galerkin operator.  This operator is traditionally organized in a block fashion where the outer sparsity pattern is dictated by the stochastic discretization while each block follows the sparsity pattern of the PDE spatial discretization.  To improve performance on architectures such as GPUs, Stokhos commutes this layout to obtain an outer sparsity pattern dictated by the spatial discretization where each block follows the stochastic discretization and leverages Kokkos hierarchical parallelism to implement calculations corresponding to the stochastic discretization in a scalable and efficient manner.}

\Comment{Page 16, subsec 4.4, par 3 -- ``.. Tpetra solver stack as a template...'' What does this mean? Also at the end of the paragraph ``union of the information'' What information? Totally unclear.}

\Response{We revised this paragraph to better explain how Stokhos interacts with Tpetra through its scalar type template parameter by adding:}

\Changes{This type can then be used within the Tpetra solver stack as an argument for its scalar type template parameter, allowing instantiation of Tpetra data structures on the ensemble type to propagate ensembles through supported linear solvers and preconditioners.}

\Response{as well as be more explicit on what is meant by ``information'' by revising the sentence as:}

\Changes{for example, GMRES for ensemble types needs to monitor the convergence of the individual samples in order to decide when to terminate based on whether every sample within the ensemble has satisfied the chosen stopping criteria.}

\Comment{Page 16, subsec 4.5, par 1 -- ``performing parameter continuation'' Isn’t this carried out by the LOCA package? I am confused. Is this another abstraction interface?}

\Response{Yes, Piro provides a unified interface and utilizes LOCA for parameter continuation, as explained in the third bullet of section 4.5. We have revised the section to clarify this point.}

\Comment{Page 16, subsec 4.5, 1st bullet point -- The text here sounds way too high-level and vague. It reads to me as there is very little practical information.}

\Response{We have revised the text to include more specific information. However, we prefer not to dedicate excessive space to the description of the Piro package.}

\Comment{Page 17, subsec 4.5, 1st bullet point -- Clarify what LOCA does and what Piro does. As it is I find it quite confusing. As a general comment, I also do not understand why interfaces and solvers have to be packed in separate packages. It increases the confusion.}

\Response{Piro integrates several packages and solvers into a single interface, which cannot be confined to individual packages. However, each package maintains its own interfaces.}

\Response{We have added the following text to clarify the roles of LOCA and Piro:}

\Changes{The core capabilities for continuation and bifurcation analysis are provided through the package LOCA. The Piro driver integrates these functionalities with linear, nonlinear, and time integration solvers, enhancing flexibility and usability.}

\noindent\textbf{References}

\Comment{Be careful to unify the convention adopted with authors name. compare for example the first
author of [17], [18] with [19], [20].}

\Response{Fixed.}

%%%%%%%%%%%%%%%%%%%%%%%%%%%%%%%%%%%%%%%%%%%%%%%%%%%%%%%
%% Reviewer 3
%%%%%%%%%%%%%%%%%%%%%%%%%%%%%%%%%%%%%%%%%%%%%%%%%%%%%%%
\pagebreak
\Review{3}

%%%%%%%%
\Comment{Comments:
This manuscript presents a comprehensive update on the Trilinos software
framework, detailing its evolution, current architecture, core packages, and
integration strategies for modern high-performance computing (HPC) environments.}

\Comment{The manuscript is not original in the traditional scientific sense (e.g.,
presenting new theorems or experimental results), but rather offers a
state-of-the-art survey/update of a large-scale software ecosystem. Nonetheless,
the documentation of such a mature, evolving scientific infrastructure is highly
valuable for the computational science community.}

\Comment{The major contributions of the manuscript include an overview of the design
decisions behind the recent shift towards Tpetra/Kokkos, a detailed description
of the core packages, and an outline of the integration strategies for
heterogeneous computing.}

\Comment{The technical content is generally robust, but due to the review nature, many
methodological details (e.g., numerical performance, limitations) are glossed
over.}

\Comment{Rather than being a glorified `ls -lah` of the Trilinos repository, the
manuscript could have benefited from a more critical analysis of the design
choices, performance trade-offs, and usability aspects of the framework,
including at least some quantitative benchmarks for key packages, justifying the
major shift that Trilinos has undergone in recent years.}

\Comment{An overview of quantitative performance metrics for key packages (e.g., MueLu,
FROSch, ROL), showing performance on various architectures, scaling trends, and
comparing to earlier versions of Trilinos would enrich the manuscript
significantly.}

\Response{We agree that a discussion of performance results could make the manuscript more interesting for readers with an interest in particular algorithms of packages of Trilinos.
  However, we do not think that this is feasible as it would significantly increase the scope.
  A lot more algorithmic detail would be required in order to provide a reasonable explanation of what platforms and scaling experiments are being run.
  Moreover, selecting a subset of problems and packages to illustrate a representative snapshot of ``performance of Trilinos'' seems extremely daunting given the breadth of algorithms.
  The current manuscript is already well beyond the usual page limit of the journal.
}

\Comment{Details on versioning, release engineering (TriBITS), and continuous integration
are well-covered, but could be expanded providing a table of build/test matrix
configurations, showing supported platforms (e.g., CUDA, ROCm, SYCL). Clarify
how to contribute performance-tuned external Kokkos-based packages into
Trilinos.}

\Comment{While this is a software-focused paper, several experiments could significantly
strengthen its utility, strong and weak scaling plots of FROSch, MueLu, and ROL
on at least two architectures (e.g., AMD+ROCm vs. NVIDIA+CUDA), to validate
performance portability claims.}

\Response{Please see our previous response.}

\Comment{The manuscript is a high-quality overview of a critical scientific
infrastructure. With the addition of quantitative benchmarks and minor editorial
improvements, it will serve as a lasting reference for both developers and
end-users in the HPC community.}

\end{document}
